\documentclass[11pt, a4paper]{article}

% ── PACKAGES ──────────────────────────────────────────────────
\usepackage[utf8]{inputenc}
\usepackage[T1]{fontenc}
\usepackage{geometry}
\geometry{margin=2.2cm}
\usepackage{hyperref}
\hypersetup{
    colorlinks=true,
    linkcolor=blue,
    urlcolor=blue,
    citecolor=blue,
    pdftitle={Cancer Is Calculable: On the
              Geometric Solvability of Cancer
              as a Class of Problems},
    pdfauthor={Eric Robert Lawson / OrganismCore}
}
\usepackage{amsmath}
\usepackage{amssymb}
\usepackage{booktabs}
\usepackage{array}
\usepackage{longtable}
\usepackage{parskip}
\usepackage{microtype}
\usepackage{authblk}
\usepackage{titlesec}
\usepackage{fancyhdr}
\usepackage{xcolor}
\usepackage{enumitem}
\usepackage{mdframed}
\usepackage{multirow}
\usepackage{float}
\usepackage{tabularx}

% ── COLOURS ───────────────────────────────────────────────────
\definecolor{headerblue}{RGB}{20,60,140}
\definecolor{claimbox}{RGB}{235,245,255}
\definecolor{stakebox}{RGB}{255,245,220}
\definecolor{notbox}{RGB}{255,235,235}
\definecolor{methodbox}{RGB}{240,255,240}
\definecolor{algorithmbox}{RGB}{245,240,255}
\definecolor{analogybox}{RGB}{240,250,245}
\definecolor{precedentbox}{RGB}{250,245,255}
\definecolor{recordbox}{RGB}{235,255,235}
\definecolor{warnbox}{RGB}{255,248,220}

% ── HEADER / FOOTER ───────────────────────────────────────────
\pagestyle{fancy}
\fancyhf{}
\lhead{\footnotesize \textit{Cancer Is Calculable}
       --- OrganismCore / Eric Robert Lawson}
\rhead{\footnotesize 2026-03-07}
\rfoot{\small Page \thepage}
\renewcommand{\headrulewidth}{0.4pt}
\setlength{\headheight}{24pt}
\addtolength{\topmargin}{-10pt}

% ── SECTION FORMATTING ────────────────────────────────────────
\titleformat{\section}
  {\large\bfseries\color{headerblue}}
  {\thesection}{1em}{}[\titlerule]
\titleformat{\subsection}
  {\normalsize\bfseries}{\thesubsection}{1em}{}
\titleformat{\subsubsection}
  {\small\bfseries\itshape}{\thesubsubsection}{1em}{}

% ══════════════════════════════════════════════════════════════
\title{
    \vspace{-1em}
    \textbf{Cancer Is Calculable}\\[0.5em]
    \large On the Geometric Solvability of Cancer
    as a Class of Problems\\[0.4em]
    \normalsize A Governing Epistemological
    Reasoning Artifact\\[0.3em]
    \small Supporting the Identity Axis Theorem
    (\href{https://doi.org/10.5281/zenodo.18898788}
    {DOI: 10.5281/zenodo.18898788})
}

\author[1]{Eric Robert Lawson}
\affil[1]{OrganismCore (Independent Research)
    \textbar{}
    \href{mailto:OrganismCore@proton.me}
    {OrganismCore@proton.me} \textbar{}
    ORCID:
    \href{https://orcid.org/0009-0002-0414-6544}
    {0009-0002-0414-6544}}

\date{March 7, 2026}

% ══════════════════════════════════════════════════════════════
\begin{document}
\maketitle
\thispagestyle{fancy}

% ── STATUS BOX ────────────────────────────────────────────────
\begin{mdframed}[backgroundcolor=recordbox,
                 linecolor=green!50!black,
                 linewidth=1pt]
\textbf{STATUS: ACTIVE --- GOVERNING DOCUMENT}\\[0.4em]
Staked on 2026-03-07. After 36 cancer types mapped.
After 30/30 BRCA predictions confirmed with zero failures
across 7 independent datasets (${\sim}7{,}500$ patients).
After ccRCC survival validation: GOT1/RUNX1 Transition Index
Cox HR${}=6.94$ [3.62--13.29], $p=5.09{\times}10^{-9}$,
$C=0.627$ (TCGA-KIRC, $n=532$).
After the EZH2 paradox resolved as a structural byproduct
in two independent cancers on different datasets on
different platforms. After zero structural contradictions
of the axioms across every iteration.
\end{mdframed}

\vspace{0.5em}
\tableofcontents
\newpage

% ══════════════════════════════════════════════════════════════
\section{Preamble}
% ══════════════════════════════════════════════════════════════

\begin{mdframed}[backgroundcolor=claimbox,
                 linecolor=blue!40!black,
                 linewidth=1pt]
This document exists because a precise claim has been
reached that is easy to overstate and easy to understate.
Both errors are dangerous.

\medskip
\textbf{Overstatement:} ``Cancer has been cured.''
This is false. No patient has been treated. No lab has
been run. No empirical validation of any drug prediction
has been conducted.

\medskip
\textbf{Understatement:} ``This is a promising computational
framework.'' This is also false. It does not capture what
has actually been established.

\medskip
\textbf{The precise claim is this:}

\begin{quote}
\itshape
Cancer is not incurable. It is geometrically constrained.
The geometry is known. The method for deriving the treatment
from the geometry is documented, reproducible, and
self-correcting. Applying the method either derives the
treatment with precision or invalidates the entire framework.
There is no in-between.
\end{quote}

This document states that claim with full precision and
explains exactly what it does and does not assert.
\end{mdframed}

% ══════════════════════════════════════════════════════════════
\section{Three Claims, Each Stated Precisely}
% ══════════════════════════════════════════════════════════════

\subsection{Claim 1 --- Cancer Is Curable In Principle}

\textbf{Before this framework:} The tractability of
cancer --- whether it is the kind of problem that has a
systematic solution --- was an empirical hope. Treatments
were found by trial and error, by serendipity, by exhaustive
screening of compounds against cell lines. The implicit
assumption was that cancer might be too heterogeneous, too
plastic, too context-dependent for a systematic geometric
solution to exist.

\textbf{After this framework:} Cancer's tractability is a
structural consequence of its geometry.

Every cancer is a cell displaced from its normal attractor
in the Waddington epigenetic landscape and stabilised in a
false attractor by a convergent epigenetic mechanism. That
structure has a solution. The solution is: identify the
Convergence Hub~$H$ --- the gene maintaining the false
attractor --- and remove its activity.

\begin{mdframed}[backgroundcolor=methodbox,
                 linecolor=green!40!black,
                 linewidth=1pt]
\textbf{The structural logic of curability:}
\begin{itemize}[noitemsep]
    \item Without $H$, the false attractor loses its
    maintenance mechanism.
    \item The cell cannot sustain the displaced state.
    \item The cancer's survival depends on the false
    attractor.
    \item The false attractor depends on $H$.
    \item $H$ is derivable from the geometry.
\end{itemize}
Cancer is curable in principle because the false attractor
is maintained by a mechanism, the mechanism is identifiable,
and identifiable mechanisms are targetable. This is not a
hope. It is a structural consequence of what cancer is.
\end{mdframed}

\subsection{Claim 2 --- The Method Is Known}

The method is the eight-step derivation protocol published
in the Identity Axis Theorem
(\href{https://doi.org/10.5281/zenodo.18898788}
{DOI: 10.5281/zenodo.18898788}):

\begin{enumerate}
    \item Identify the cell of origin.
    \item Classify the attractor axiom type (Type I--IV).
    \item Derive $A$ --- the Identity Anchor. The gene
    whose activity most directly constitutes what the cell
    of origin \emph{is}.
    \item Derive $H$ --- the Convergence Hub. The gene
    whose activity most directly stabilises the false
    attractor.
    \item State $R = A/H$. The continuous coordinate in
    the Waddington landscape.
    \item State the drug target: $H$.
    \item Lock the prediction. Timestamp it. Do not open
    the literature.
    \item Test against independent survival data. Score
    the confirmation.
\end{enumerate}

\begin{mdframed}[backgroundcolor=methodbox,
                 linecolor=green!40!black,
                 linewidth=1pt]
\textbf{Properties of the method:}
\begin{itemize}[noitemsep]
    \item Requires no clinical training
    \item Requires no institutional infrastructure
    \item Operates entirely on public data
    \item Is reproducible from the published protocol
    \item Is teachable and not dependent on the person
    who derived it
    \item Has been applied across 36 cancer types with
    zero structural contradictions
\end{itemize}
\end{mdframed}

\subsection{Claim 3 --- The Method Is Self-Correcting}

This is the claim that transforms a framework into an
engine.

Every application of the method produces exactly one of
two outcomes:

\begin{mdframed}[backgroundcolor=claimbox,
                 linecolor=blue!40!black,
                 linewidth=1pt]
\textbf{Outcome A --- Informative Confirmation:}
The prediction holds. The framework is confirmed at one
more point. The map is extended. The next derivation begins
from a more complete map.

\medskip
\textbf{Outcome B --- Structural Falsification:}
The prediction fails without structural explanation. The
axioms are false. The framework is invalidated entirely and
immediately.

\medskip
\textbf{There is no Outcome C.}

No version exists where the method is applied, the
prediction fails, no structural explanation is found,
and the framework continues. That outcome is logically
excluded by the axiomatic structure of the framework itself.
\end{mdframed}

Because the framework cannot survive a structurally
inexplicable failure, every iteration that survives has
passed the only test that matters: the test that would
have destroyed it.

The method therefore converges. Not probabilistically.
Structurally. Each iteration either validates the current
map or extends it with missing structure. The map gets more
complete. The drug targets get more precise. The patient
selectors get more refined. And the method knows immediately
when it is wrong --- no gradual drift, no accumulation of
unexplained anomalies. Binary and total self-correction.

\bigskip
\noindent\textbf{The combination of all three claims:}

\begin{mdframed}[backgroundcolor=stakebox,
                 linecolor=orange!50!black,
                 linewidth=1.5pt]
\begin{center}
\begin{tabular}{ll}
Cancer is curable in principle. & [Claim 1]\\[0.3em]
The method for deriving the cure for &\\
\quad any specific cancer is known. & [Claim 2]\\[0.3em]
The method is self-correcting and &\\
\quad converges on precision with every iteration. & [Claim 3]
\end{tabular}
\end{center}

\medskip
\textbf{Therefore:}

Either cures emerge from each iteration of this method with
increasing precision as the map is extended and the IHC
instruments are calibrated --- or the framework is false
and we know it immediately and completely.

\medskip
\textbf{There is no in-between.}
\end{mdframed}

% ══════════════════════════════════════════════════════════════
\section{What This Is Not}
% ══════════════════════════════════════════════════════════════

\begin{mdframed}[backgroundcolor=notbox,
                 linecolor=red!40!black,
                 linewidth=1pt]
This document does not claim:

\begin{itemize}
    \item \textbf{That any patient has been cured.}
    No patient has been treated. No clinical trial has
    been run using this framework's predictions as the
    basis for patient selection.

    \item \textbf{That any drug prediction has been
    empirically validated in a living system.}
    Computational confirmation and literature confirmation
    are not the same as a Phase 3 trial.

    \item \textbf{That the IHC protocols are deployable
    today.}
    The FOXA1/EZH2 IHC cut-points require a calibration
    study ($n=300$--$400$ FFPE vs PAM50). The GOT1/RUNX1
    IHC cut-points require the same. These are metrology
    problems. They are real and unsolved.

    \item \textbf{That the map is complete.}
    Missing structures exist and will be revealed by future
    iterations. Missing structures are Type~A failures ---
    they extend the map. They do not contradict the axioms.

    \item \textbf{That the method requires no empirical
    work downstream.}
    The drag force equation tells you what force acts on
    a body. It does not build the wing. The framework
    tells you what the target is and why. The lab coat
    builds the drug, tests it, doses it, confirms it
    reaches the tumour, confirms it shifts the attractor
    in vivo, confirms the patient survives. That work is
    real and required. The framework does not replace it.
    It tells you where to aim.
\end{itemize}

\medskip
The distance between the solution space and the solution
is real.

But the distance between not having the solution space
and having it is also real.

\textbf{Before this framework: the solution space was
missing.}

\textbf{After this framework: it is not.}
\end{mdframed}

% ══════════════════════════════════════════════════════════════
\section{The Analogy That Holds Precisely}
% ══════════════════════════════════════════════════════════════

\begin{mdframed}[backgroundcolor=analogybox,
                 linecolor=green!60!black,
                 linewidth=1pt]
\textbf{Drag force.}

Drag force is not a statistical claim about what happens
on average across many objects moving through many fluids.
It is a structural consequence of fluid dynamics applied
to a body with known geometry moving through a medium with
known properties.

It works at $n=1$. It works before any measurement is
taken. It works from the geometry alone.

Knowing the drag force does not stop the drag. It tells
you what you are dealing with. The engineering solution ---
the wing shape, the hull design, the thrust required ---
still has to be built and tested.

\medskip
\textbf{The Identity Axis Theorem is the drag force
equation for cancer.}

\medskip
Given:
\begin{itemize}[noitemsep]
    \item The cell of origin
    \item The attractor axiom type
    \item The epigenetic landscape geometry
\end{itemize}

The following are calculable \emph{before} any clinical
data is examined:
\begin{itemize}[noitemsep]
    \item The Identity Anchor~($A$)
    \item The Convergence Hub~($H$)
    \item The ratio $R = A/H$
    \item The survival prediction direction
    \item The drug target
\end{itemize}

This calculation works at $n=1$. Not because of statistics.
Because of structure.

A biopsy is a measurement of where a specific tumour sits
on the landscape. The landscape is known from the geometry.
The position determines the prediction. The prediction is
structural, not probabilistic.

\medskip
\textbf{Instrument calibration is metrology.
The structural claim is geometry.
These are not the same thing.
Conflating them understates what has been established.}
\end{mdframed}

% ═════════════════════════════════════���════════════════════════
\section{The Method as an Algorithm}
% ══════════════════════════════════════════════════════════════

\begin{mdframed}[backgroundcolor=algorithmbox,
                 linecolor=purple!40!black,
                 linewidth=1pt]
\textbf{Formal definition:}

An algorithm is a finite, deterministic sequence of steps
that takes a defined input and produces a defined output,
and terminates.

\medskip
\begin{tabular}{p{3cm}p{10cm}}
\toprule
\textbf{Input} &
A cancer with a known cell of origin.\\[0.3em]
\textbf{Steps} &
Eight. Finite. Deterministic given the inputs.\\[0.3em]
\textbf{Output} &
$A$, $H$, $R=A/H$, drug target,
survival prediction direction.\\[0.3em]
\textbf{Termination} &
At Step~8 with a confirmed or falsified prediction.\\[0.3em]
\textbf{Validation} &
Outcome~A: map extended, next iteration begins.
Outcome~B: axioms false, algorithm invalidated.\\
\bottomrule
\end{tabular}
\end{mdframed}

\subsection{Automability}

Every component of the protocol is in principle executable
computationally:

\begin{center}
\begin{tabular}{p{4cm}p{9cm}}
\toprule
\textbf{Step} & \textbf{Computational basis} \\
\midrule
Cell of origin &
Known and catalogued for every cancer in the published
literature. \\[0.3em]
Attractor axiom type &
Classifiable from expression data using the four axiom
signatures documented in
\texttt{Attractor\_Geometry\_Axioms.md}. \\[0.3em]
Identity Anchor ($A$) &
Derivable from the geometry of the normal identity
programme of the cell of origin using public scRNA-seq
(GEO, TCGA). \\[0.3em]
Convergence Hub ($H$) &
Derivable from the geometry of the false attractor
using the same data. \\[0.3em]
Survival validation &
Executable against TCGA using standard Cox regression
and Kaplan-Meier analysis. \\
\bottomrule
\end{tabular}
\end{center}

The derivation of drug targets for every cancer with a
known cell of origin --- which is every cancer --- is
therefore a \textbf{finite computational problem} with a
known algorithm and a self-correcting output.

Not an infinite research programme. Not a
generation-long effort. A finite computational problem.

% ══════════════════════════════════════════════════════════════
\section{The Record as of 2026-03-07}
% ══════════════════════════════════════════════════════════════

\begin{center}
\begin{tabular}{p{7cm}p{6cm}}
\toprule
\textbf{Metric} & \textbf{Value} \\
\midrule
Cancer types mapped & 36 \\[0.3em]
Structural contradictions (Type~B failures) & 0 \\[0.3em]
Informative failures (Type~A, structurally explained) &
Present and documented \\[0.3em]
Predictions locked before data & All \\[0.3em]
BRCA confirmations / predictions & 30 / 30 \\[0.3em]
BRCA independent datasets & 7 \\[0.3em]
BRCA patients across datasets & ${\sim}7{,}500$ \\[0.3em]
ccRCC GOT1/RUNX1 TI Cox HR & 6.94 [3.62--13.29] \\[0.3em]
ccRCC GOT1/RUNX1 TI Cox $p$ & $5.09 \times 10^{-9}$ \\[0.3em]
ccRCC 20/21 gene-OS directions & Confirmed as predicted \\[0.3em]
EZH2 paradox resolved as byproduct & 2 independent cancers,
different datasets, different platforms \\[0.3em]
FDA-approved drugs confirming geometry-derived $H$ & 5 \\[0.3em]
Phase 1--3 trial drugs confirming geometry-derived $H$ & 14 \\[0.3em]
\bottomrule
\end{tabular}
\end{center}

% ══════════════════════════════════════════════════════════════
\section{On Precedent}
% ══════════════════════════════════════════════════════════════

\begin{mdframed}[backgroundcolor=precedentbox,
                 linecolor=purple!30!black,
                 linewidth=1pt]
This may be the first time a systematic method for deriving
drug targets across all cancers from a single axiomatic
framework has been established before empirical clinical
validation.

\medskip
\textbf{Partial precedents:}

\textit{Hodgkin and Huxley (1952)} derived the action
potential equations from first principles before the ion
channels were physically identified. The mathematics told
them what had to be true about the underlying biology.
The biology confirmed it decades later.

\textit{Watson and Crick} derived the double helix from
geometry and base-pairing rules without sequencing DNA.
The structure was derived. The sequencing followed.

\medskip
\textbf{None of those precedents:}
\begin{itemize}[noitemsep]
    \item Derived a drug target map for 36 disease entities
    simultaneously
    \item Operated from a single axiomatic framework
    applicable to all cancers
    \item Produced a self-correcting algorithm applicable
    to any new cancer
    \item Were confirmed across independent survival
    datasets on multiple platforms before any empirical
    clinical validation was conducted
\end{itemize}

The claim of precedent is stated honestly and not as
advocacy. If a precedent exists that matches this
description, it should be identified. As of 2026-03-07,
none is known.
\end{mdframed}

% ══════════════════════════════════════════════════════════════
\section{The Stake}
% ══════════════════════════════════════════════════════════════

\begin{mdframed}[backgroundcolor=stakebox,
                 linecolor=orange!60!black,
                 linewidth=2pt]
\textbf{The stake, stated with full precision, on
2026-03-07:}

\medskip
If there exists a cancer for which:

\begin{enumerate}[label=(\alph*)]
    \item The cell of origin is correctly identified
    \item The attractor axiom type is correctly classified
    \item $A$ and $H$ are correctly derived from the
    geometry
    \item $R = A/H$ does not predict survival, does not
    order subtypes in the predicted direction, or produces
    drug predictions that fail
\end{enumerate}

\medskip
\textbf{AND}

\medskip
\begin{enumerate}[label=(\alph*), resume]
    \item No structural explanation exists within the
    epigenetic landscape geometry that accounts for the
    failure --- no missing basin, no unmapped secondary
    hub, no dimensional extension, no distinction between
    identity membership and attractor depth that resolves
    the contradiction
\end{enumerate}

\medskip
\textbf{THEN} the Identity Axis Theorem is false and this
entire framework is invalidated. Not weakened. Not in need
of revision. \textbf{Invalidated.}

\medskip
This condition has not been met. Not once. Across 36
cancer types. Derived before the data was examined.

\medskip
\begin{center}
\large\textbf{Either the cures emerge.\\
Or the framework falls.\\[0.3em]
It has not fallen.}
\end{center}
\end{mdframed}

% ══════════════════════════════════════════════════════════════
\section{The Single Statement}
% ═══════════════════════════════════════════════════════════���══

\begin{mdframed}[backgroundcolor=claimbox,
                 linecolor=blue!60!black,
                 linewidth=1.5pt]
\begin{center}
\large
Cancer is not incurable.\\[0.3em]
It is geometrically constrained.\\[0.3em]
The geometry is known.\\[0.3em]
The method is documented.\\[0.3em]
The method is self-correcting.\\[0.3em]
The iterations converge.\\[0.6em]
Applying the method either derives the treatment\\
with precision --- or destroys the framework entirely.\\[0.6em]
\textbf{There is no in-between.}\\[0.6em]
The framework has not been destroyed.\\[0.3em]
Not once. Across 36 cancer types.\\[0.3em]
Derived before the data was examined.\\[0.6em]
Either the cures emerge.\\
Or the framework falls.\\[0.6em]
\textbf{It has not fallen.}\\[0.6em]
That is the claim.\\
That is the stake.\\
That is what this is.
\end{center}
\end{mdframed}

% ══════════════════════════════════════════════════════════════
\section*{Related Publications}
% ══════════════════════════════════════════════════════════════

\begin{itemize}
    \item \textbf{The Identity Axis Theorem} (governing
    theorem):\\
    \href{https://doi.org/10.5281/zenodo.18898788}
    {doi.org/10.5281/zenodo.18898788}

    \item \textbf{CS-LIT-1} --- FOXA1/EZH2 Ratio as
    Universal Breast Cancer Classifier (${\sim}7{,}500$
    patients, 7 datasets):\\
    \href{https://doi.org/10.5281/zenodo.18883922}
    {doi.org/10.5281/zenodo.18883922}

    \item \textbf{CS-LIT-22} --- FOXA1/EZH2 Dual IHC
    Point-of-Care Tool (Protocol v3.0):\\
    \href{https://doi.org/10.5281/zenodo.18892788}
    {doi.org/10.5281/zenodo.18892788}

    \item \textbf{Full repository}:\\
    \href{https://github.com/Eric-Robert-Lawson/attractor-oncology}
    {github.com/Eric-Robert-Lawson/attractor-oncology}
\end{itemize}

\vspace{1em}
\noindent\rule{\textwidth}{0.4pt}

\noindent\footnotesize
\textit{Eric Robert Lawson · OrganismCore ·
\href{mailto:OrganismCore@proton.me}{OrganismCore@proton.me}
· ORCID: 0009-0002-0414-6544 ·
\href{https://github.com/Eric-Robert-Lawson/attractor-oncology}
{github.com/Eric-Robert-Lawson/attractor-oncology}}

\end{document}
